% BATHY_PDE_APPROACH.tex
%
% Bathymetry as a coupled radiative transfer inverse problem: rationale, physics and an
% implementation specification for photic.
%
% Build: pdflatex BATHY_PDE_APPROACH.tex (three times, for the cross-references).

\documentclass[11pt]{article}
\usepackage{amsmath,amssymb,amsthm,mathtools}
\usepackage[margin=2.5cm]{geometry}
\usepackage{booktabs}
\usepackage{longtable}
\usepackage{xcolor}
\usepackage{hyperref}
\definecolor{dark-blue}{rgb}{0,0,0.45}
\hypersetup{colorlinks, linkcolor={dark-blue}, citecolor={blue}, urlcolor={blue}}

\numberwithin{equation}{section}

\newcommand{\bxi}{\boldsymbol{\xi}}
\newcommand{\bx}{\mathbf{x}}
\newcommand{\rTOA}{\rho^{\mathrm{TOA}}}
\newcommand{\robs}{\rho^{\mathrm{obs}}}
\newcommand{\dd}{\,\mathrm{d}}
\newcommand{\code}[1]{\texttt{#1}}

\title{Bathymetry as a Coupled Radiative Transfer Inverse Problem:\\
Rationale, Physics and Implementation Specification for \textit{photic}}
\author{Design specification}
\date{16 October 2026}

\begin{document}
\maketitle

\begin{abstract}
\textit{photic} inverts atmospherically corrected reflectance with Lee's semi-analytical
model, one pixel neighbourhood at a time. ACOLITE does the atmospheric correction beforehand,
for each scene separately. This document specifies a replacement for both steps, built in
stages.

The forward model becomes the radiative transfer equation (RTE), the integro-differential
transport equation that HydroLight solves. It covers the atmosphere, the sea surface, the water
column and the bottom, and it is evaluated over each band's spectral response. The RTE is never
solved inside the fit. Its exact solution is precomputed as a lookup table over a few
dimensionless variables, using a discrete-ordinate solver written for \textit{photic}. The fit
interpolates the table, which gives the speed of a table with an interpolation error that is
measured and kept below the sensor noise. The atmosphere is handled the same way, from the
6SV-based tables that ACOLITE distributes.

The inversion keeps \textit{photic}'s structure: a bounded Levenberg--Marquardt fit over a pixel
window and all scenes at once, with depth and bottom shared across scenes and water properties
per scene. In the later stages each scene's aerosol joins the fitted parameters. Part~I gives
the rationale and the physics. Part~II is the specification: modules, data formats, commands,
numerical methods, milestones, tests and acceptance criteria.
\end{abstract}

\tableofcontents

%============================================================================================
\part{Rationale and physics}

%--------------------------------------------------------------------------------------------
\section{Does this already exist?}\label{sec:prior}

We found no published method that does all of the following at once:
\begin{enumerate}
\item solves the full RTE, not a parametric approximation, for the coupled
atmosphere--water--bottom system;
\item retrieves bottom depth;
\item inverts several scenes jointly, with depth and bottom shared across them.
\end{enumerate}
Each part exists separately.

\begin{itemize}
\item \textbf{Exact water-column RT.} HydroLight solves the RTE in the water column by invariant
imbedding, for given inherent optical properties (IOPs), depth and bottom reflectance
\cite{mobley1994,oceanopticsbook}. OSOAA \cite{chami2015} and CAO-DISORT \cite{jin1994} solve
the coupled atmosphere--ocean problem. All three are forward models.
\item \textbf{Coupled atmosphere--ocean inversion.} Joint retrievals of aerosol and water
properties exist for open and coastal ocean colour. Examples are the multi-angle polarimetric
retrieval developed for PACE \cite{gao2019} and POLYMER \cite{steinmetz2011}. None models a
bottom.
\item \textbf{Bayesian shallow-water inversion.} Thompson et al.\ \cite{thompson2017} solve for
the most probable depth, seafloor reflectance and water properties. Their atmospheric
correction is a separate step run beforehand \cite{thompson2018}.
\item \textbf{Lookup tables of RT solutions.} Spectrum matching against tables of HydroLight
spectra over discrete water, bottom and depth combinations is established \cite{mobley2005}.
\item \textbf{Multi-temporal bathymetry.} Published multi-image satellite methods mostly
composite the images, then fit an empirical or machine-learning model. \textit{photic}'s
temporal extension \cite{blake2020} is the joint physical case, but it uses Lee's model on
corrected reflectance.
\end{itemize}

%--------------------------------------------------------------------------------------------
\section{Why change the forward model}\label{sec:why}

Validation of \textit{photic} found the following (\code{validation/EXPERIMENTS.md};
\code{REVIEW.md}, \S2d):
\begin{enumerate}
\item \textbf{False depths fit as well as true ones.} With one bottom type selected by fit, deep
water at the Gulf entrance gets false shallow depths: RMSE 10.8\,m, bias $-6.8$\,m. Their
misfit is as low as that of the correct solution.
\item \textbf{Errors the fit cannot see.} The tables were sampled at band centres. Averaged over
the Sentinel-2 B2 response, pure-water absorption is 32\% higher. In the synthetic tests that
gives depth errors of $-17\%$ to $+28\%$ at a reduced $\chi \le 0.15$.
\item \textbf{The correction is fixed before the inversion.} Changing ACOLITE's aerosol
retrieval (tiled, or with adjacency correction) moves tier 4 by more than the noise floor, in
opposite directions at different sites. The inversion cannot push back on an aerosol estimate
that its own fit rejects.
\item \textbf{Lee's model approximates the RTE.} It is a fit to HydroLight runs. With four
visible bands, a structural error of a few per cent becomes metres of depth.
\end{enumerate}
Joint processing of several scenes holds constraints that per-scene correction throws away.
The seabed is the same in every scene up to tide, so an aerosol estimate that makes one scene's
bottom signal consistent with the others is more likely to be right.

%--------------------------------------------------------------------------------------------
\section{The forward problem}\label{sec:forward}

\subsection{The coupled RTE}

$\bx$ is horizontal position and $z$ height (positive upwards, mean sea surface at $z=0$).
Scene $t$ has tide $\eta_t$, so its water column is $-D_t < z < 0$ with
\begin{equation}
D_t(\bx) = H(\bx) + \eta_t(\bx),
\end{equation}
where $H$ is the depth below the common datum (\textit{photic}'s \code{HTIDE} convention). For
each wavelength $\lambda$ and direction $\bxi$ ($\mu = \bxi\cdot\hat{\mathbf z}$), the radiance
satisfies
\begin{equation}\label{eq:rte3d}
\bxi\cdot\nabla L + c\,L = \int_{S^2}\beta(\bxi'\!\to\bxi)\,L(\bxi')\dd\Omega',
\end{equation}
in the atmosphere ($0<z<z_T$) and in the water. $c = a+b$ is the beam attenuation and $\beta$
the volume scattering function. The boundary conditions are:
\begin{align}
&z=z_T,\ \mu<0: && L = F_0\,\delta(\bxi-\bxi_0), \label{eq:bc-toa}\\
&z=0: && \begin{pmatrix} L^{\uparrow}(0^+)\\ L^{\downarrow}(0^-)\end{pmatrix}
  = \begin{pmatrix} \mathcal{R}_{aw} & \mathcal{T}_{wa}\\ \mathcal{T}_{aw} & \mathcal{R}_{wa}\end{pmatrix}
    \begin{pmatrix} L^{\downarrow}(0^+)\\ L^{\uparrow}(0^-)\end{pmatrix}, \label{eq:bc-surf}\\
&z=-D_t,\ \mu>0: && L = \frac{\rho_b}{\pi}\int_{\mu'<0}|\mu'|\,L(\bxi')\dd\Omega'. \label{eq:bc-bot}
\end{align}
The surface operators are Fresnel reflection and transmission, averaged over the Cox--Munk
distribution of wave slopes \cite{coxmunk1954} for a rough surface. The bottom is Lambertian,
with $\rho_b(\bx,\lambda) = B(\bx)\,\hat\rho_{k(\bx)}(\lambda)$: one shape $\hat\rho_k$ from the
bottom library per pixel (\textit{photic}'s single-type rule) and a brightness $B$.

\subsection{Water IOPs}

These are \textit{photic}'s present IOP model, with the scattering coefficient and phase
function made explicit:
\begin{align}
a &= a_w + (a_0 + a_1\ln P)_+\,P + G\,e^{-S_g(\lambda-440)}, \label{eq:iop-a}\\
b_{bp} &= X\,(440/\lambda)^{Y}, \qquad b_p = b_{bp}/B_p, \qquad b_w = 2\,b_{bw}, \label{eq:iop-b}\\
\beta &= b_w\,\tilde\beta_w + b_p\,\tilde\beta_{\mathrm{FF}}(\,\cdot\,;B_p). \label{eq:iop-beta}
\end{align}
\begin{itemize}
\item $\tilde\beta_w(\psi) \propto 1 + 0.835\cos^2\psi$ is the pure-water phase function.
\item $\tilde\beta_{\mathrm{FF}}$ is the Fournier--Forand phase function \cite{fournier1994} with
backscatter fraction $B_p$, fixed by default at 0.0183, the Petzold average.
\item The tables $a_w$, $b_{bw}$, $a_0$ and $a_1$ come from \code{model/data}.
\item $(\cdot)_+$ is the existing clamp at zero.
\end{itemize}

\subsection{The measurement}

Band $j$ of scene $t$ is the response-weighted average
\begin{equation}\label{eq:obs}
R_{rs,j} = \frac{\int S_j(\lambda)\,F_0(\lambda)\,R_{rs}(\lambda)\dd\lambda}{\int S_j(\lambda)\,F_0(\lambda)\dd\lambda}
\end{equation}
for water-leaving reflectance (stages M1 and M2). At top of atmosphere (stage M3) the same
average is applied to $\rTOA$. $S_j$ comes from \code{model/data/srf}; $F_0$ is the
extraterrestrial solar irradiance (new data file, \S\ref{sec:data}).

\subsection{A hierarchy of approximations}

\begin{description}
\item[Column (levels 0--1).] The water and bottom are a plane-parallel column at each pixel,
exact for horizontally uniform water and bottom. The atmosphere is plane-parallel. At level
1, neighbouring pixels are coupled through the atmospheric adjacency kernel
\cite{tanre1981,vermote1997}.
\item[3D (level 2).] Monte Carlo or SHDOM \cite{evans1998}, used only to check the column
approximation at reef edges. It is never used inside the fit.
\end{description}

%--------------------------------------------------------------------------------------------
\section{Lookup tables from the PDE solution}\label{sec:lut}

The RTE is solved offline and stored as a table that the fit interpolates. The table's only
error is interpolation, which can be measured and reduced. The structural error of Lee's
fitted model cannot.

\subsection{Why the table is small}

\begin{enumerate}
\item \textbf{Similarity.} For a homogeneous column, the monochromatic RTE depends on the IOPs and
depth only through three quantities:
\begin{equation}\label{eq:similarity}
\omega = \frac{b}{c}, \qquad \tau = c\,D, \qquad f_w = \frac{b_w}{b}.
\end{equation}
$\omega$ is the single-scattering albedo and $\tau$ the optical depth. $f_w$ is the share of
molecular scattering, which sets the mixed phase function \eqref{eq:iop-beta} for a fixed
$B_p$. Wavelength is not a dimension: \eqref{eq:iop-a}--\eqref{eq:iop-b} map each wavelength
to its own $(\omega,\tau,f_w)$. So the same table serves every band and every response
sub-sample. The interface depends on wavelength only through the refractive index of water,
which is held at 1.34.

\item \textbf{The bottom in closed form.} For a Lambertian bottom, the remote-sensing reflectance
just above the surface, in any direction, is exactly
\begin{equation}\label{eq:bottom-closed}
R_{rs}(\rho_b) = R_0 + \frac{A\,\rho_b}{1 - C\,\rho_b}.
\end{equation}
\begin{itemize}
\item $R_0$ is the reflectance over a black bottom.
\item $A = E_{b0}\,L_e$, where $E_{b0}$ is the irradiance reaching a black bottom and $L_e$ the
radiance leaving the surface per unit irradiance emitted isotropically by the bottom.
\item $C$ is the fraction of that emitted irradiance which the column and surface return to the
bottom.
\end{itemize}
None of these depends on $\rho_b$. Two solutions with the same operator determine them: one
over a black bottom, and one with an isotropically emitting bottom and no sky.

\item \textbf{One solve gives every direction.} A discrete-ordinate solution gives the radiance
in all directions, so view angle and azimuth add storage but no solves.
\end{enumerate}

\subsection{Illumination}

The table is built for two illuminations of the surface: the direct sun at $\theta_s$, and a
uniform sky of unit irradiance. A scene band's reflectance is
\begin{equation}\label{eq:illum}
R_{rs} = (1-s)\,R_{rs}^{\mathrm{dir}} + s\,R_{rs}^{\mathrm{sky}},
\end{equation}
where $s$ is the diffuse fraction of the downwelling irradiance. $s$ comes from the atmosphere
tables at the scene's aerosol optical depth, or from a Rayleigh-only estimate when none is
given. The uniform sky is an approximation. Its effect on near-nadir $R_{rs}$ is small and is
measured in test T6 (\S\ref{sec:tests}).

%--------------------------------------------------------------------------------------------
\section{The inverse problem}\label{sec:inverse}

\subsection{Unknowns}

\begin{center}
\begin{tabular}{@{}llll@{}}
\toprule
Symbol & Meaning & Shared across scenes? & Stage\\
\midrule
$H$ & depth below datum, per region of the window & yes & M1\\
$k$, $B$ & bottom type (discrete), brightness per region & yes & M1\\
$\ln P, \ln G, \ln X$ & water IOPs per scene (window-constant) & no & M1\\
$\Delta_t$ & glint offset per scene (\code{GLINT\_OFFSET}) & no & M1\\
$\tau_{a,t}$ & aerosol optical depth at 550\,nm per scene & no & M3\\
$m_t$ & aerosol model per scene (discrete) & no & M3\\
$\delta\eta_t$ & tide-datum offset per scene & no & M4\\
\bottomrule
\end{tabular}
\end{center}

\subsection{Objective}

For each window (\textit{photic}'s pixel neighbourhood), the fit minimises
\begin{equation}\label{eq:objective}
\chi^2 = \sum_{t,j,r}\frac{\big(R_{t,j}^{\mathrm{model}}(\theta) - R_{t,j}^{\mathrm{obs}}\big)_r^2}{\sigma_{t,j}^2}
 + \sum_i \frac{(\theta_i - \theta_i^{\mathrm{prior}})^2}{\sigma_i^2},
\end{equation}
by bounded Levenberg--Marquardt, as \code{samodel\_lm} does now. The prior rows are the existing
continuity, coherence and depth-prior rows, plus the aerosol prior of stage M3. At stage M3,
$R$ is the TOA reflectance and the bands include NIR and SWIR.

\subsection{Gradients}

At stages M1--M4 the forward model is local to the window, so the Jacobian follows from the
chain rule through the table interpolant (\S\ref{sec:jacobian}). No adjoint is needed. The
adjoint RTE becomes necessary only for the optional global solve (M6), in which the
atmosphere and adjacency couple all pixels. The design for that is in
Appendix~\ref{app:adjoint}.

\subsection{Discrete choices}

The bottom type and aerosol model are chosen by block-coordinate descent:
\begin{enumerate}
\item fit the continuous parameters for each candidate;
\item keep the candidate with the lowest $\chi^2$ (\code{SELECT\_BY\_FIT}).
\end{enumerate}
The optically deep alternative is the existing test, with $R_0(\tau\to\infty)$ from the
table.

%============================================================================================
\part{Implementation specification}

%--------------------------------------------------------------------------------------------
\section{Scope and milestones}\label{sec:milestones}

Each milestone is a separate branch and change. Its forward model is switchable, and the
existing model stays the default until the acceptance rule of \code{validation/README.md}
(tier 3 and tier 4 depth metrics improve, with no regressions) says otherwise.

\begin{longtable}{@{}p{1.1cm}p{5.6cm}p{7.6cm}@{}}
\toprule
& Deliverable & Done when\\
\midrule
\endhead
M0 & Band averaging of the existing model (\code{SCENE ... SRF}); spectral data in files. &
Done on branch \code{data-inputs}; benchmarked; adopted as the acceptance rule decides.\\
M1 & \code{rtwater.c}: discrete-ordinate solver for the water column. &
Passes tests T1--T4 (\S\ref{sec:tests}).\\
M2 & \code{rttable.c}: per-scene tables, interpolation and derivatives. RT forward model in
\code{samodel.c} behind \code{SET FORWARD RT}; ACOLITE $R_{rs}$ input. &
T5--T8 pass; \code{make test-rt} passes; tier 2 RT-vs-OSOAA error $<1\%$; tier 3 and 4
benchmarked against M0.\\
M3 & \code{atmos.c}: reader for ACOLITE's LUTs; TOA forward model; aerosol per scene in the fit;
NIR/SWIR bands. &
T9--T11 pass; tier 4 run from L1 TOA reflectance, compared with ACOLITE input.\\
M4 & Tide-datum offsets per scene; optically deep test from the table. &
Tier 4, including Gulf entrance false depths, benchmarked.\\
M5 & Adjacency (level 1). & Reef-edge 3D check (level 2); tier 4.\\
M6 & Optional global solve with the adjoint (Appendix~\ref{app:adjoint}). & Research; not
scheduled.\\
\bottomrule
\end{longtable}

%--------------------------------------------------------------------------------------------
\section{Software architecture}\label{sec:arch}

New modules follow the conventions in \code{CLAUDE.md}: C17, 2-space indentation in new files,
\code{//} comments in sentence case, British spelling, globals declared in \code{common.h} and
defined in \code{common.c}. OpenMP is not used.

\begin{longtable}{@{}p{2.6cm}p{11.8cm}@{}}
\toprule
File & Responsibility\\
\midrule
\endhead
\code{rtwater.c/.h} & Discrete-ordinate solution of one homogeneous water layer under a Fresnel
interface, over a black or isotropically emitting Lambertian bottom (\S\ref{sec:solver}). Pure
numerics, with no \textit{photic} state.\\
\code{rttable.c/.h} & Table axes, build, cache file I/O, validation against direct solves, and
interpolation with derivatives (\S\ref{sec:table}).\\
\code{atmos.c/.h} (M3) & ACOLITE LUT reader; band-averaged path reflectance, transmittances,
spherical albedo and diffuse fraction as functions of $\tau_a$, with derivatives
(\S\ref{sec:atmos}).\\
\code{spectra.c} & Gains the solar irradiance $F_0$ and the band sub-sampling used by
\eqref{eq:obs}.\\
\code{samodel.c} & A forward-model switch in \code{samodel\_Rrs} and the Jacobian code
(\S\ref{sec:samodel}); aerosol parameters (M3).\\
\code{bam.c} & Commands and \code{SCENE} keywords (\S\ref{sec:commands}).\\
\code{model/tests/} & \code{rt\_check.c} (T1--T8), and \code{synth\_gen} in RT mode for
\code{make test-rt}.\\
\code{validation/} & \code{tier2\_osoaa.py} gains an RT-mode comparison; \code{rt\_table\_check.py}
compares tables with OSOAA (T4).\\
\bottomrule
\end{longtable}

\paragraph{Dependencies.} \code{rtwater.c} needs a real symmetric or nonsymmetric eigensolver for
matrices up to $64\times64$ and an LU solver. Write them in C (Householder reduction and QR
iteration, as in DISORT's \code{ASYMTX} \cite{stamnes1988}), or use LAPACK if the build adds
it. The default is in-tree code, so that the dependencies stay NetCDF and PGPLOT.

%--------------------------------------------------------------------------------------------
\section{The water-column solver (\code{rtwater.c})}\label{sec:solver}

\subsection{Problem}

The solver handles a homogeneous slab of optical thickness $\tau$ and single-scattering albedo
$\omega$, with a phase function given by its Legendre moments $\chi_\ell$, $\ell = 0..2N-1$. It
lies under an interface with relative refractive index $n_w = 1.34$, on a Lambertian bottom.
It returns, per unit downwelling irradiance just above the surface:
\begin{itemize}
\item the remote-sensing reflectance just above the surface at the requested view direction;
\item the irradiance reaching the bottom;
\end{itemize}
for the three source cases: direct sun, uniform sky, and emitting bottom with no sky.

\subsection{Discretisation}

\begin{itemize}
\item \textbf{Streams.} Use $N$ streams per hemisphere in the water (default $N=24$, at most 32).
They are split into a refracted cone ($0$ to the critical angle, $\arcsin(1/n_w) \approx
48.3^\circ$) and a total-internal-reflection region, each with its own double-Gauss quadrature,
as CAO-DISORT does \cite{jin1994}. This keeps the interface's discontinuity at the critical
angle on a quadrature boundary.
\item \textbf{Azimuth.} Expand in Fourier modes $m = 0..M$, with $M = 2N-1$. The Lambertian
bottom couples only $m=0$. The direct beam drives every mode. The sky drives only $m=0$.
\item \textbf{Phase function.} The Fournier--Forand moments are computed once per $B_p$ by
Gauss--Legendre integration of the analytic phase function, with 2000 or more nodes
concentrated at small angles. Use $\delta$-M truncation \cite{wiscombe1977}, and correct the
single scattering into the view direction with the full phase function, as in Nakajima and
Tanaka's TMS method \cite{nakajima1988}. The FF phase function is strongly forward-peaked, so
without truncation $N$ would have to be in the hundreds.
\item \textbf{Interface.} At M1--M2 the interface is flat: Fresnel reflection and transmission
are diagonal in the quadrature. Radiance in the refracted cone is mapped through Snell's law
with the $n^2$ radiance factor, and total internal reflection applies beyond the critical
angle. In the stage after M3, Cox--Munk rough-surface operators are integrated numerically per
Fourier mode, for a wind speed $w$.
\item \textbf{Solution.} In each mode, a homogeneous layer has the eigen solution of
\cite{stamnes1988}, with the reduced $N\times N$ eigenproblem. The direct beam adds a
particular solution. The boundary conditions form a $2N\times2N$ linear system:
\begin{itemize}
\item at the top, downward radiance is the transmitted sky plus $\mathcal{R}_{wa}$ applied to
the upward radiance;
\item at the bottom, upward radiance is $\rho/\pi$ times the downward irradiance for $m=0$, and
zero for $m>0$; for the emission case, a unit isotropic source.
\end{itemize}
The sources share one factorisation of each mode's system.
\item \textbf{Output.} Evaluate the upwelling radiance just below the surface at the view
direction refracted into the water from the discrete-ordinate solution, using the source
function integration of \cite{stamnes1988}, then transmit it through the interface. Return
$R_0$, $E_{b0}$, $L_e$ and $C$ for \eqref{eq:bottom-closed}, for the direct and sky sources.
\end{itemize}

\subsection{API}

\begin{verbatim}
typedef struct {
  int n_streams;            // per hemisphere, default 24
  double n_water;           // 1.34
  double bp;                // FF backscatter fraction, default 0.0183
  double wind;              // 0: flat surface (M1-M2)
} rt_settings;

typedef struct {
  double R0, A, C;          // R_rs = R0 + A rho/(1 - C rho), per E_d(0+)
} rt_coeffs;

// Prepare moments, quadrature and interface operators; 0 on success.
int rt_init(const rt_settings *s, rt_workspace **w);
// Solve one node (omega, tau, f_w); geometry in air, degrees.
int rt_solve(rt_workspace *w, double omega, double tau, double f_w,
             double theta_sun, double theta_view, double dphi,
             rt_coeffs *direct, rt_coeffs *sky);
void rt_free(rt_workspace *w);
\end{verbatim}
The optically deep limit, $\tau\to\infty$, is solved directly: there is no bottom boundary,
and the upward radiance is the layer's reflection eigenvector solution. It gives $R_0^\infty$
with $A = 0$.

\subsection{Cost}

The target is $\le 1$\,ms per node on one core, with $N=24$ and all modes. Most of the cost
is the eigenproblems and the boundary-condition factorisation. They are shared by all three
sources and by every view direction.

%--------------------------------------------------------------------------------------------
\section{Tables (\code{rttable.c})}\label{sec:table}

\subsection{Axes}

\begin{center}
\begin{tabular}{@{}llll@{}}
\toprule
Axis & Range & Transform & Nodes (default)\\
\midrule
$\omega$ & 0.02--0.995 & $u = \ln(1-\omega)$ & 32\\
$\tau$ & $10^{-3}$--$\tau_{\max}$ (40), plus $\infty$ & $v = \ln\tau$ & 48 + deep\\
$f_w$ & 0--1 & linear & 7\\
\bottomrule
\end{tabular}
\end{center}
Geometry is fixed per table: one table per scene, for that scene's $(\theta_s, \theta_v,
\Delta\phi)$ in air, each rounded to $0.1^\circ$ for the cache key. Every node stores $(R_0, A,
C)$ for the direct and sky sources, 6 doubles, so a table is $32\times49\times7\times6\times8$
bytes $\approx$ 0.5\,MB.

\subsection{Build}

Build all nodes with \code{rt\_solve}; 11\,000 nodes take about 10--20\,s per scene at the
target cost. Then validate (below). Tables are cached as NetCDF in \code{\$PHOTIC\_CACHE}
(default \code{\$HOME/.cache/photic/rt}),
keyed by a hash of:
\begin{itemize}
\item the geometry;
\item the \code{rt\_settings};
\item the axes;
\item the solver version string.
\end{itemize}
A scene whose table is in the cache costs nothing.

\paragraph{Cache file format.} The file has the dimensions \code{omega}, \code{tau},
\code{fw}, \code{source} (2: direct, sky) and \code{coeff} (3: R0, A, C), and the variable
\code{table(source, coeff, fw, tau, omega)} as a double. It holds the axis variables and the
following global attributes:
\begin{itemize}
\item \code{theta\_sun}, \code{theta\_view} and \code{delta\_phi};
\item \code{n\_streams}, \code{n\_water}, \code{bp} and \code{wind};
\item \code{solver\_version};
\item the validation statistics: \code{max\_rel\_error}, \code{p95\_rel\_error} and
\code{n\_checks}.
\end{itemize}

\subsection{Interpolation}\label{sec:interp}

Use tensor-product cubic Hermite interpolation in $(u,v,f_w)$. The node derivatives come from
centred finite differences of the node values in the transformed coordinates (one-sided at the
ends). This makes the interpolant $C^1$ and exact for cubics. Interpolate the coefficients
$R_0$, $A$ and $C$, never $R_{rs}$ itself, so \eqref{eq:bottom-closed} stays exact in $\rho_b$.
Each evaluation returns the value and its three partial derivatives.

\subsection{Validation of a table}

After a build, solve 200 random off-node points directly, drawn uniformly in the transformed
coordinates and $\rho_b\in[0,1]$. Compare $R_{rs}$ for $s = 0.2$. Record the maximum and 95th
percentile relative errors in the file. Then:
\begin{itemize}
\item if the 95th percentile exceeds 0.3\% or the maximum exceeds 1\%, double the nodes along
the axis with the largest error contribution, and rebuild;
\item after two refinements, print a warning and continue.
\end{itemize}
The threshold is about a tenth of tier 4's noise level ($R_{rs}$ noise about $5\times10^{-4}$
on signals about $5\times10^{-3}$).

\subsection{Derivatives for the fit}\label{sec:jacobian}

For one band sub-sample, with $D = H + \eta_t$,
\begin{align*}
\omega &= b/c, & \tau &= cD, & f_w &= b_w/b, & c &= a+b, & b &= b_w + b_p.
\end{align*}
The partial derivatives of the similarity variables with respect to the IOPs are
\begin{align}
\frac{\partial\omega}{\partial a} &= -\frac{b}{c^2}, & \frac{\partial\omega}{\partial b_p} &= \frac{a}{c^2},
& \frac{\partial\tau}{\partial a} &= \frac{\partial\tau}{\partial b_p} = D, & \frac{\partial\tau}{\partial H} &= c,\\
\frac{\partial f_w}{\partial b_p} &= -\frac{b_w}{b^2}, & \frac{\partial f_w}{\partial a} &= 0. \nonumber
\end{align}
The IOP derivatives with respect to $\ln P$, $\ln G$, $\ln X$ and $Y$ are those of
\code{water\_iops}, with $b_p = b_{bp}/B_p$. For the bottom,
\begin{equation}
\frac{\partial R_{rs}}{\partial\rho_b} = \frac{A}{(1-C\rho_b)^2}, \qquad
\frac{\partial\rho_b}{\partial B} = \hat\rho_k.
\end{equation}
Band averaging \eqref{eq:obs} is linear, so the Jacobian rows are the same weighted sums of
the sub-sample derivatives. All derivatives are those of the interpolant, so the
Gauss--Newton model and the depth $\sigma$ (\code{DEPTH\_SIGMA}) are consistent with the
misfit the fit evaluates.

%--------------------------------------------------------------------------------------------
\section{Changes to \code{samodel.c}}\label{sec:samodel}

\begin{enumerate}
\item \textbf{Forward model switch.} Add a global \code{forward\_model}, set by \code{SET
FORWARD LEE|RT}, default LEE. Declare it in \code{common.h} and define it in \code{common.c}.
In RT mode:
\begin{itemize}
\item \code{samodel()} gets or builds each scene's table before the pixel loop, and prints
the table's validation statistics;
\item \code{samodel\_Rrs} computes each band from its sub-samples, \eqref{eq:obs} and
\eqref{eq:illum}, through the table;
\item it fills \code{rrs\_dp} with $R_0^{\infty}$ and \code{rrs\_bottom} with the bottom term,
which the optically deep test and the index of optical depth use.
\end{itemize}
\item \textbf{Sub-samples.} For a scene with \code{SRF}, the sub-samples are 5\,nm steps across
the band's response, with weights $S_jF_0\Delta\lambda$, precomputed per scene and band. The
IOP tables are interpolated at each sub-sample. Without \code{SRF}, there is one sub-sample at
the band wavelength.
\item \textbf{Jacobians.} \code{region\_jacobian} and \code{water\_jacobian} get RT branches built
from \S\ref{sec:jacobian}. Until they exist, RT mode uses the existing finite-difference path
(\code{ANALYTIC\_JACOBIAN 0}), which is correct but slower.
\item \textbf{Removed in RT mode.} Lee's coefficients (\code{LEE\_*}), the $D_u$ path elongations,
\code{INTERREFLECTION} (now exact in \eqref{eq:bottom-closed}) and the conversion
$0.5\,r_{rs}/(1-1.5\,r_{rs})$ are all subsumed by the table.
\item \textbf{Kept.} The parameter vector, bounds, prior rows, \code{GLINT\_OFFSET},
\code{Y\_MODEL}, single-type selection, profile search, LUT warm starts, depth $\sigma$ and
outputs.
\item \textbf{The cache of IOPs.} \code{iop\_cache} is extended with the sub-sample $(\omega, \tau/D,
f_w)$ per scene, so that only $\tau$ is recomputed when $H$ changes.
\end{enumerate}

%--------------------------------------------------------------------------------------------
\section{The atmosphere (\code{atmos.c}, M3)}\label{sec:atmos}

\subsection{Data}

ACOLITE distributes 6SV-based LUTs per sensor, already averaged over each band's response
(\code{ACOLITE-LUT-202110}). There is one NetCDF file per aerosol model (MOD1 continental,
MOD2 maritime) and per pressure (500, 750, 1013 and 1100\,hPa). Each band variable has the
dimensions \code{(par, azi, thv, ths, wnd, tau)}. The \code{par} attribute lists the
parameters, including \code{romix} (path reflectance), \code{utott} and \code{dtott} (total
up and down transmittances), \code{astot} (spherical albedo) and the direct components. The
axes are relative azimuth 0--180$^\circ$, view and sun zenith 0--80$^\circ$, and
$\tau_{550}$ 0.001--5. They are not part of \textit{photic}'s repository: ACOLITE's licence
applies and they total 1.3\,GB. They are read from the user's ACOLITE installation.

\subsection{TOA model}

\begin{equation}\label{eq:toa}
\rTOA_j = T_{g,j}\Big[\rho_{\mathrm{path},j}(\tau_a)
  + \frac{T^{\downarrow}_j(\tau_a)\,T^{\uparrow}_j(\tau_a)\,\pi R_{rs,j}}{1 - S_j(\tau_a)\,\pi R_{rs,j}}
  + T^{\mathrm{dir}}_j\,\Delta_t\Big].
\end{equation}
\begin{itemize}
\item $\rho_{\mathrm{path}}$, $T^\downarrow$, $T^\uparrow$ and $S$ are interpolated from the
LUT: multilinear in angle and pressure, cubic in $\ln\tau_a$, with derivatives.
\item $T_g$ is the gas transmittance, from ACOLITE's gas LUT with the ozone and water vapour
from the L1 metadata, as ACOLITE does.
\item $R_{rs,j}$ is the band's water-leaving reflectance from the table, with the sky fraction
$s = 1 - T^{\downarrow,\mathrm{dir}}/T^{\downarrow}$ in \eqref{eq:illum}.
\item $\Delta_t$ is the existing glint offset, now in TOA units.
\end{itemize}

\subsection{In the fit}

Add $\tau_{a,t}$ as a per-scene water-block parameter, bounded to $[0.001, 1.5]$. Its prior
comes from a dark-spectrum estimate over the scene's deep and dark pixels in SWIR, which
\code{atmos.c} computes once per scene, as ACOLITE's DSF does. The prior's $\sigma$ is 0.03,
or as set. The aerosol model is a discrete choice per scene, made in that first pass by the
better SWIR fit, then fixed. The bands are the scene's visible, NIR and SWIR bands. Over water
the NIR and SWIR water-leaving terms are taken as zero, which pins $\tau_a$.

\subsection{Inputs}

Use L1 TOA reflectance grids, written by \code{validation/tier4\_prepare.py} from ACOLITE's L1R
output (\code{l1r\_export\_geotiff}). They are resampled to the site grid like the L2W bands.
Scenes get the keyword \code{TOA} (\S\ref{sec:commands}).

%--------------------------------------------------------------------------------------------
\section{Commands and data}\label{sec:commands}\label{sec:data}

\subsection{Commands}

\begin{longtable}{@{}p{6.4cm}p{8cm}@{}}
\toprule
Command & Effect\\
\midrule
\endhead
\code{SET FORWARD LEE|RT} & Forward model of \code{MODEL Lee\_Semi\_Analytical} (default LEE).\\
\code{SET RTSTREAMS n} & Streams per hemisphere (default 24).\\
\code{SET RTBP x} & Fournier--Forand backscatter fraction (default 0.0183).\\
\code{SET RTCACHE dir} & Table cache directory.\\
\code{COMPUTE RTTABLE SCENE s} & Build (or load) scene $s$'s table and print its validation.\\
\code{PRINT RTTABLE SCENE s} & Print a table's settings and statistics.\\
\code{SCENE ... SRF f [AOT x]} & M2: $\tau_{550}$ for the sky fraction (default: from \code{ATMOSLUT}
if given, else Rayleigh only).\\
\code{SCENE ... SRF f TOA ATMOSLUT dir} & M3: the bands are TOA reflectance; ACOLITE LUT directory
\code{[PRESSURE hPa] [AEROSOL MOD1|MOD2|FIT] [OZONE cm] [WATERVAPOUR cm]}.\\
\bottomrule
\end{longtable}
\code{SCENE} keywords stay positional and optional, after \code{SRF}, in the order shown. RT mode
requires \code{SRF} on every scene. Without it, photic prints a warning and uses one
sub-sample at the band wavelength.

\subsection{New data files}

\begin{itemize}
\item \code{model/data/f0\_thuillier\_2003.csv}: extraterrestrial solar irradiance, 1\,nm,
350--1000\,nm, from Thuillier et al.\ \cite{thuillier2003}, for the weights in \eqref{eq:obs}.
\item \code{model/data/srf/}: add the NIR and SWIR bands for M3 (S2 B8, B8A, B11, B12; OLI B5,
B6, B7).
\item No RT tables are committed: they are built per scene and cached. The optional global
table (Appendix~\ref{app:global}) would be a downloadable data file.
\end{itemize}

%--------------------------------------------------------------------------------------------
\section{Tests}\label{sec:tests}

\subsection{Unit and solver tests (\code{model/tests/rt\_check.c}, \code{make test-rt})}

\begin{longtable}{@{}p{0.8cm}p{6.4cm}p{7.2cm}@{}}
\toprule
& Test & Pass\\
\midrule
\endhead
T1 & Energy: $\omega=1$, white bottom; flux in equals flux out. & Relative error $<10^{-6}$.\\
T2 & Thin layer: $\tau=10^{-4}$, compared with analytic single scattering plus Fresnel
reflection of the bottom. & $<0.1\%$.\\
T3 & Convergence in $N$: 16, 24, 32 and 48 streams. & $R_{rs}$ changes by $<0.2\%$ from 24
to 48.\\
T4 & Against OSOAA (scalar build), on the tier 2 grid of water types, depths, bottoms and sun
angles. & Median $<0.5\%$, 95th percentile $<1.5\%$, for $\omega\le0.95$.\\
T5 & Closed form \eqref{eq:bottom-closed}: direct solves at $\rho_b = 0, 0.3, 0.7, 1$ against the
coefficients. & $<10^{-6}$.\\
T6 & Uniform-sky approximation: a uniform sky against a Rayleigh sky, from the atmosphere
tables, for $\theta_s = 20$--$60^\circ$. & Report; expected $<1\%$ at near-nadir view.\\
T7 & Table interpolation: the 200-point validation of every table. & 95th percentile $<0.3\%$,
maximum $<1\%$.\\
T8 & Jacobian: interpolant derivatives against centred differences of the interpolant. &
$<10^{-5}$ relative.\\
T9 & ACOLITE LUT reader: \code{romix}, \code{dtott} etc.\ at nodes, against direct reads of the
NetCDF file. & Exact at nodes.\\
T10 & TOA model: ACOLITE's own $\rho_w$ for a scene, put through \eqref{eq:toa} at ACOLITE's
$\tau_a$, against the L1R TOA. & Agreement within ACOLITE's conventions; report the
difference.\\
T11 & Aerosol retrieval on synthetic TOA scenes (OSOAA with aerosol, PHOTIC\_TIER3\_SKY). &
$\tau_a$ within 0.02.\\
\bottomrule
\end{longtable}

\code{make test-rt} also runs \code{synth\_gen} in RT mode: synthetic scenes from the RT
forward model, inverted by \code{MODEL Lee\_Semi\_Analytical} with \code{SET FORWARD RT}. Its
outputs are compared with a committed baseline, as \code{make test} does now. The existing
\code{make test} is unchanged.

\subsection{Benchmark}

\begin{itemize}
\item \textbf{Tier 2.} RT mode against OSOAA on the existing cases. This is T4 at scale, and
replaces the tier's role of measuring Lee's structural error.
\item \textbf{Tier 3.} The \code{\_insitu} scenarios, with bottoms from held-out measured
spectra. They are the main synthetic test. A new water truth, OSOAA at 5\,nm sub-samples and
then band-averaged, makes tier 3 independent of the band-centre tables. It is built once and
cached.
\item \textbf{Tier 4.} Exmouth on Sentinel-2 and Landsat 8 (Tantabiddi, Bundegi, Marina, Gulf
entrance); later, the Arabian Gulf site with ICESat-2 ATL24 truth. M3 runs from L1 TOA
reflectance.
\item \textbf{Acceptance.} Use the existing rule: tier 3 and tier 4 depth metrics must improve,
with no regressions beyond the noise floor of about $\pm0.5$ points or $\pm0.3$\,m RMSE at one
site. Run time is reported. The target is at most $2\times$ the present model once each
scene's table is cached.
\end{itemize}

%--------------------------------------------------------------------------------------------
\section{Risks and open questions}\label{sec:risks}

\begin{itemize}
\item \textbf{The bottom library.} An exact forward model cannot correct a bottom shape that does
not match the seabed. In the \code{\_insitu} tier 3 scenarios, fitting a single type per pixel
against heterogeneous measured bottoms is much worse than fitting mixtures (RMSE 5--6\,m
against about 2\,m in clear water). The bottom parameterisation needs its own decision,
independent of the RT work.
\item \textbf{Aerosol--water trade-off.} With four visible bands, $\tau_a$ and $X$ can trade off.
The defences are the SWIR bands, the shared seabed across scenes and the prior from the dark
spectrum.
\item \textbf{Sky illumination.} The uniform sky approximation is quantified by T6. If it matters,
add an atmosphere-derived sky radiance distribution as a third source.
\item \textbf{Raman scattering and fluorescence} are not in the elastic RTE. Raman adds a few per
cent in clear blue water. A correction table can be added later.
\item \textbf{Stratified water} and non-Lambertian bottoms are outside the table.
\item \textbf{Surface roughness} is flat in M1--M2. Cox--Munk operators come with the TOA model.
\item \textbf{Licences.} ACOLITE's LUTs are read in place, not redistributed. The in situ
spectra are CC BY 3.0 and their attribution is in the data file.
\end{itemize}

%--------------------------------------------------------------------------------------------
\appendix
\section{A global table (optional)}\label{app:global}

The per-scene table can be replaced by one global table with the sun angle (12 nodes) and view
direction (8 zenith nodes, with azimuthal Fourier coefficients) as further axes. That is about
$10^5$ solves, or minutes on a workstation, and a file of some tens to a few hundred megabytes.
It is worth doing only if per-scene builds become a bottleneck, for example with many short
scenes.

\section{Global solve with the adjoint (M6, optional)}\label{app:adjoint}

When the atmosphere is a smooth field coupled across all pixels, or adjacency is modelled,
the inversion becomes one PDE-constrained problem over the whole site:
\begin{equation}
\min_{u}\ \tfrac12\sum_{t,j}\Big\|\frac{\rTOA_{t,j}[u]-\robs_{t,j}}{\sigma_{t,j}}\Big\|^2
 + \alpha_H\,\mathrm{TV}_\epsilon(H) + \alpha_w\sum_t\|\nabla\theta^t_w\|^2 + \alpha_a\sum_t\|\Delta\theta^t_a\|^2.
\end{equation}
Write the RTE and its boundary conditions as $\mathcal{A}(u)L = q$. The adjoint radiance $L^*$
solves $\mathcal{A}^*L^* = \sum_j\mathcal{M}_j^*\,(\rTOA_j-\robs_j)/\sigma_j^2$. The gradient
is then $\partial\mathcal{J}/\partial u_i = -\langle L^*,\partial_{u_i}\mathcal{A}\,L\rangle$,
whatever the number of unknowns \cite{box1988,hinze2009}.

For depth, map the column to $s = z/D\in(-1,0)$. $D$ then becomes a coefficient,
$\mu\,\partial_sL = D(-cL + \int\beta L)$, and the gradient is
\begin{equation}
\frac{\partial\mathcal{J}}{\partial H} = \sum_t\int_{-1}^{0}\!\!\int L_t^{*}\Big(-cL_t + \int\beta L_t\Big)\dd\Omega\dd s.
\end{equation}
Matrix-free Gauss--Newton uses tangent-linear products for $Jv$ and adjoint products for
$J^{\mathsf T}w$. With the table forward model, only the atmospheric and adjacency couplings
need adjoints; the column part is the table's own derivative.

%--------------------------------------------------------------------------------------------
\begin{thebibliography}{99}

\bibitem{blake2020} S.~Blake, \textit{photic}: a physics-based shallow water bathymetry model
for multispectral imagery, arXiv:2002.02298 (2020).

\bibitem{box1988} M.~A. Box, S.~A.~W. Gerstl, C.~Simmer, Application of the adjoint
formulation to the calculation of atmospheric radiative effects, \textit{Beitr.\ Phys.\
Atmos.}\ 61 (1988) 303--311.

\bibitem{chami2015} M.~Chami, B.~Lafrance, B.~Fougnie, J.~Chowdhary, T.~Harmel, F.~Waquet,
OSOAA: a vector radiative transfer model of coupled atmosphere--ocean system for a rough sea
surface, \textit{Opt.\ Express} 23 (2015) 27829--27852.

\bibitem{coxmunk1954} C.~Cox, W.~Munk, Measurement of the roughness of the sea surface from
photographs of the sun's glitter, \textit{J.\ Opt.\ Soc.\ Am.}\ 44 (1954) 838--850.

\bibitem{evans1998} K.~F. Evans, The spherical harmonics discrete ordinate method for
three-dimensional atmospheric radiative transfer, \textit{J.\ Atmos.\ Sci.}\ 55 (1998)
429--446.

\bibitem{fournier1994} G.~R. Fournier, J.~L. Forand, Analytic phase function for ocean water,
\textit{Proc.\ SPIE} 2258, Ocean Optics XII (1994) 194--201.

\bibitem{gao2019} M.~Gao, P.-W. Zhai, B.~A. Franz, et al., Inversion of multiangular
polarimetric measurements over open and coastal ocean waters: a joint retrieval algorithm for
aerosol and water-leaving radiance properties, \textit{Atmos.\ Meas.\ Tech.}\ 12 (2019)
3921--3941.

\bibitem{hinze2009} M.~Hinze, R.~Pinnau, M.~Ulbrich, S.~Ulbrich, \textit{Optimization with PDE
Constraints}, Springer (2009).

\bibitem{jin1994} Z.~Jin, K.~Stamnes, Radiative transfer in nonuniformly refracting layered
media: atmosphere--ocean system, \textit{Appl.\ Opt.}\ 33 (1994) 431--442.

\bibitem{mobley1994} C.~D. Mobley, \textit{Light and Water: Radiative Transfer in Natural
Waters}, Academic Press (1994).

\bibitem{mobley2005} C.~D. Mobley, L.~K. Sundman, C.~O. Davis, et al., Interpretation of
hyperspectral remote-sensing imagery by spectrum matching and look-up tables, \textit{Appl.\
Opt.}\ 44 (2005) 3576--3592.

\bibitem{nakajima1988} T.~Nakajima, M.~Tanaka, Algorithms for radiative intensity
calculations in moderately thick atmospheres using a truncation approximation, \textit{J.\
Quant.\ Spectrosc.\ Radiat.\ Transfer} 40 (1988) 51--69.

\bibitem{oceanopticsbook} C.~D. Mobley (ed.), Ocean Optics Web Book,
\url{https://www.oceanopticsbook.info/view/introduction/overview}.

\bibitem{stamnes1988} K.~Stamnes, S.-C. Tsay, W.~Wiscombe, K.~Jayaweera, Numerically stable
algorithm for discrete-ordinate-method radiative transfer in multiple scattering and emitting
layered media, \textit{Appl.\ Opt.}\ 27 (1988) 2502--2509.

\bibitem{steinmetz2011} F.~Steinmetz, P.-Y. Deschamps, D.~Ramon, Atmospheric correction in
presence of sun glint: application to MERIS, \textit{Opt.\ Express} 19 (2011) 9783--9800.

\bibitem{tanre1981} D.~Tanr\'e, M.~Herman, P.-Y. Deschamps, Influence of the background
contribution upon space measurements of ground reflectance, \textit{Appl.\ Opt.}\ 20 (1981)
3676--3684.

\bibitem{thompson2017} D.~R. Thompson, E.~J. Hochberg, G.~P. Asner, et al., Airborne mapping
of benthic reflectance spectra with Bayesian linear mixtures, \textit{Remote Sens.\ Environ.}\
200 (2017) 18--30.

\bibitem{thompson2018} D.~R. Thompson, V.~Natraj, R.~O. Green, et al., Optimal estimation for
imaging spectrometer atmospheric correction, \textit{Remote Sens.\ Environ.}\ 216 (2018)
355--373.

\bibitem{thuillier2003} G.~Thuillier, M.~Herse, D.~Labs, et al., The solar spectral
irradiance from 200 to 2400 nm as measured by the SOLSPEC spectrometer from the ATLAS and
EURECA missions, \textit{Solar Phys.}\ 214 (2003) 1--22.

\bibitem{vermote1997} E.~F. Vermote, D.~Tanr\'e, J.-L. Deuz\'e, M.~Herman, J.-J. Morcrette,
Second Simulation of the Satellite Signal in the Solar Spectrum, 6S: an overview, \textit{IEEE
Trans.\ Geosci.\ Remote Sens.}\ 35 (1997) 675--686.

\bibitem{wiscombe1977} W.~J. Wiscombe, The delta-M method: rapid yet accurate radiative flux
calculations for strongly asymmetric phase functions, \textit{J.\ Atmos.\ Sci.}\ 34 (1977)
1408--1422.

\end{thebibliography}

\end{document}
